\section*{Capítulo \ref{ch:b3:origin-of-elements} --- De onde vêm os elementos}
% Values from the Book 5 scratch calculator calc/ch30.py.

\begin{solution}{exo:b3:origin-of-elements:1}
(a)~$\mathrm{Si/Fe} = 10^{7.51 - 7.46} = 10^{0.05} = 1.12$ e
$\mathrm{O/Fe} = 10^{8.69 - 7.46} = 10^{1.23} = 17.0$: cerca de um
átomo de silício e dezessete de oxigênio por átomo de ferro. (b)~$\mathrm{Fe/Au} =
10^{7.46 - 0.91} = 10^{6.55} = \num{3.55e6}$ e $\mathrm{Pb/Au} =
10^{1.95 - 0.91} = 10^{1.04} = 11.0$: além do pico do ferro, as abundâncias
caem milhões de vezes, mas o chumbo, no seu próprio pico, é
onze vezes mais abundante que o ouro.
\end{solution}

\begin{solution}{exo:b3:origin-of-elements:2}
(a)~$10^{7.46 - 12} = \num{2.884e-5}$, ou seja, \num{28.8} átomos de ferro por
milhão de átomos de hidrogênio. (b)~Cem vezes menos: \num{0.288} por
milhão, um átomo de ferro para cada $1/\num{2.884e-7} = \num{3.47e6}$ átomos de
hidrogênio. (c)~$[\mathrm{Fe/H}] = \log_{10}(\num{e-7}/\num{2.884e-5}) =
-2.46$.
\end{solution}

\begin{solution}{exo:b3:origin-of-elements:3}
$^{124}\mathrm{Xe}$ e $^{126}\mathrm{Xe}$ situam-se no lado rico em prótons,
fora do caminho s e blindados pelos estáveis $^{124}\mathrm{Te}$ e
$^{126}\mathrm{Te}$: \omterm{def:b3:origin-of-elements:p-process}{núcleos p}. $^{128}\mathrm{Xe}$ e $^{130}\mathrm{Xe}$
estão no caminho s, mas blindados do \omterm{def:b3:origin-of-elements:s-process}{processo r} pelos estáveis
$^{128}\mathrm{Te}$ e $^{130}\mathrm{Te}$: apenas s. $^{129}\mathrm{Xe}$,
$^{131}\mathrm{Xe}$ e $^{132}\mathrm{Xe}$ estão no caminho s e não têm
isóbaro estável de $Z$ menor (o iodo tem apenas 127): ambos os processos.
$^{134}\mathrm{Xe}$ e $^{136}\mathrm{Xe}$ situam-se além de
$^{133}\mathrm{Xe}$, onde o caminho s deixa o xenônio, e não são
blindados: apenas r.
\end{solution}

\begin{solution}{exo:b3:origin-of-elements:4}
(a)~$t = d/c = 40\times\qty{3.086e22}{m}/\qty{2.998e8}{m.s^{-1}} =
\qty{4.12e15}{s}$, cerca de \num{130} milhões de anos. (b)~$\Delta v/c \le
1.7/\num{4.12e15} = \num{4.1e-16}$. Trata-se de um limite superior, pois parte ou
a totalidade do atraso pode vir da emissão dos raios gama após a
coalescência.
\end{solution}

\begin{solution}{exo:b3:origin-of-elements:5}
(a)~$295\times0.1067 = \qty{31.48}{mb}$, $832\times0.0373 =
\qty{31.03}{mb}$, $155\times0.1984 = \qty{30.75}{mb}$. (b)~A partir de
$^{122}\mathrm{Te}$: $N_{123} = 31.48/832 = \num{0.03783}$, \qty{1.4}{\percent}
acima do valor medido, e $N_{124} = 31.48/155 = \num{0.2031}$,
\qty{2.4}{\percent} acima. Abundâncias que diferem por um fator de cinco são
previstas com erro de poucos por cento apenas a partir das seções de choque.
(c)~$^{128}\mathrm{Te}$ é apenas r: o fluxo s deixa o telúrio em
$^{127}\mathrm{Te}$, que decai, de modo que nenhum fluxo s estacionário passa por
$^{128}\mathrm{Te}$ e sua abundância foi determinada pelo \omterm{def:b3:origin-of-elements:s-process}{processo r}.
\end{solution}

\begin{solution}{exo:b3:origin-of-elements:6}
(a)~$\tau_0\langle\sigma\rangle = 0.306\times155 = 47.4$, logo
$f = 1/(1 + 1/47.4) = \num{0.979}$ para $^{124}\mathrm{Te}$; para
$^{138}\mathrm{Ba}$, $0.306\times4.0 = 1.224$ e $f = 1/(1 + 0.817) =
\num{0.550}$: o núcleo mágico transmite apenas \qty{55}{\percent} do
fluxo. (b)~Os seis fatores são \num{0.959}, \num{0.550}, \num{0.908},
\num{0.771}, \num{0.972} e \num{0.915}; seu produto é \num{0.328}:
$\sigma N$ cai por um fator de \num{3.05}. (c)~A queda medida é
$27.6/7.72 = 3.57$. O modelo acerta o tamanho do degrau a menos de
\qty{15}{\percent}, quase todo ele vindo de $^{138}\mathrm{Ba}$ e
$^{140}\mathrm{Ce}$, os dois núcleos com $N = 82$.
\end{solution}

\begin{solution}{exo:b3:origin-of-elements:7}
(a)~$kT = 8\times10^3\times\qty{1.602e-19}{J} = \qty{1.282e-15}{J}$, e
$v_T = \sqrt{2\times\num{1.282e-15}/\num{1.6749e-27}} =
\qty{1.24e6}{m.s^{-1}}$. (b)~Com $\langle\sigma\rangle =
\qty{155e-31}{m^2}$, $\lambda = n_\mathrm{n}\langle\sigma\rangle v_T$
é \qty{1.92e-11}{s^{-1}} para $n_\mathrm{n} = \qty{e12}{m^{-3}}$
($10^6$ por centímetro cúbico): um tempo médio de \qty{5.2e10}{s}, cerca de
\num{1650} anos; depois \num{16.5} anos a $10^8$ e \num{0.165}~ano,
dois meses, a \qty{e10}{cm^{-3}}. (c)~$^{127}\mathrm{Te}$ decai em
horas, muito mais rápido que qualquer desses tempos de captura: ele sempre decai e
o caminho não se bifurca ali. $^{134}\mathrm{Cs}$ vive cerca de dois
anos, entre os tempos de captura a $10^8$ e \qty{e10}{cm^{-3}}: nas
altas densidades dos flashes de hélio, boa parte dele é capturada
antes de decair. É um ponto de ramificação, e as abundâncias em cada lado
dele medem a densidade de nêutrons.
\end{solution}

\begin{solution}{exo:b3:origin-of-elements:8}
(a)~$10^{6.30 - 3.14} = 1445$ vezes menos abundante que o cálcio e
$10^{4.97 - 3.14} = 68$ vezes menos abundante que o titânio. (b)~$10^{8.46 - 1.38}
= \num{1.2e7}$. (c)~O escândio mostra o efeito ímpar--par: núcleos com
número par de prótons são mais fortemente ligados e são formados pela
captura de núcleos de hélio nos estágios de queima das estrelas massivas, que os
favorecem. O berílio mostra o vale do lítio, do berílio e do boro:
estes núcleos frágeis são destruídos nas temperaturas do interior estelar,
e a queima do hélio vai diretamente do hélio ao carbono,
sem produzi-los.
\end{solution}

\begin{solution}{exo:b3:origin-of-elements:9}
(a)~Seja $\alpha$ o conjunto dos \omterm{def:b3:origin-of-elements:alpha}{elementos alfa} e $\mathrm{Fe}_\mathrm{cc}$,
$\mathrm{Fe}_\mathrm{Ia}$ o ferro dos dois tipos de supernova. No
platô, $\alpha/\mathrm{Fe}_\mathrm{cc} = 10^{P}$ vezes a
razão solar. Em $[\mathrm{Fe/H}] = 0$, $\alpha/(\mathrm{Fe}_\mathrm{cc} +
\mathrm{Fe}_\mathrm{Ia})$ é solar. Dividindo,
$\mathrm{Fe}_\mathrm{cc}/(\mathrm{Fe}_\mathrm{cc} + \mathrm{Fe}_\mathrm{Ia})
= 10^{-P}$, e a fração do Tipo~Ia é $1 - 10^{-P} = 1 - 10^{-0.4} =
\num{0.602}$. (b)~$10^{-P} = 1/3$ dá $P = \log_{10}3 = 0.48$.
(c)~As supernovas do Tipo~Ia começam após o mesmo atraso nas duas galáxias.
Quando começaram na galáxia anã, seu gás havia atingido apenas
$[\mathrm{Fe/H}] = -2$, contra $-1$ na Galáxia: ela havia transformado
seu gás em estrelas dez vezes mais lentamente.
\end{solution}

\begin{solution}{exo:b3:origin-of-elements:10}
(a)~Para o zircônio,
\[
\frac{(X/\mathrm{Mg})_\mathrm{Terra}}{(X/\mathrm{Mg})_\mathrm{CI}} =
\frac{10.5/228000}{3.96/95870} = 1.11 ;
\]
analogamente, $(2670/228000)/(5010/95870) = 0.224$ para o sódio,
\num{0.190} para o potássio e \num{0.118} para o rubídio. (b)~O zircônio, o mais refratário, permanece
próximo de sua proporção condrítica. Os três metais alcalinos moderadamente voláteis,
que condensam a \qtyrange{800}{1006}{K}, estão empobrecidos por
fatores de cinco a nove, sendo o rubídio, o mais volátil, o mais empobrecido. A
Terra silicática segue a \omterm{def:b3:origin-of-elements:volatility}{tendência de volatilidade}. (c)~O magnésio é um elemento maior,
refratário, que condensa junto com a maior parte da rocha, e é
litófilo: permaneceu no manto quando o núcleo se formou. Uma referência siderófila,
como o ferro ou o níquel, misturaria a perda para o núcleo com a
perda por volatilidade.
\end{solution}

\begin{solution}{exo:b3:origin-of-elements:11}
(a)~Quando uma massa $\mathrm{d}s$ é retida em estrelas, o vento remove
$\lambda\,\mathrm{d}s$ de gás de \omterm{def:b3:origin-of-elements:metallicity}{metalicidade} $Z$:
$\mathrm{d}M_\mathrm{g} = -(1 + \lambda)\,\mathrm{d}s$ e
$\mathrm{d}(ZM_\mathrm{g}) = y\,\mathrm{d}s - Z\,\mathrm{d}s -
\lambda Z\,\mathrm{d}s$. Expandindo a segunda e subtraindo
$Z\,\mathrm{d}M_\mathrm{g}$, obtém-se $M_\mathrm{g}\,\mathrm{d}Z =
y\,\mathrm{d}s = -\frac{y}{1 + \lambda}\,\mathrm{d}M_\mathrm{g}$, que
integrada fornece o resultado. (b)~$\ln(1/0.29) = 1.238$: $Z = 0.0155$ para
$\lambda = 0$, e a metade disso, $0.0077$, para $\lambda = 1$. (c)~Um vento
reduz o rendimento efetivo pelo fator $1 + \lambda$. Em uma galáxia de
massa pequena, as supernovas expulsam o gás com mais facilidade, $\lambda$ é grande,
e os metais são perdidos junto com o gás.
\end{solution}

\begin{solution}{exo:b3:origin-of-elements:12}
(a)~$N_\mathrm{Li}/N_\mathrm{H} = 10^{0.96 - 12} = \num{9.1e-12}$ no
Sol e $10^{3.25 - 12} = \num{1.78e-9}$ nos condritos CI. O Sol
conserva $10^{-2.29} = \num{0.0051}$ de seu lítio: destruiu
\qty{99.5}{\percent} dele. (b)~$\num{1.78e-9}/\num{1.45e-10} = 12.3$ e
$\num{1.78e-9}/\num{4.72e-10} = 3.8$. (c)~O Sistema Solar contém quase
quatro vezes mais lítio do que o Big Bang produziu: a maior parte de seu lítio
foi produzida depois, por algum outro processo. As estrelas mais antigas contêm três vezes
menos do que o Big Bang produziu: ou destruíram parte dele, como
o Sol fez, ou a previsão deixa escapar algo; as duas comparações
sozinhas não dizem qual dessas hipóteses é a correta.
\end{solution}

\begin{solution}{pb:b3:origin-of-elements:1}

\textbf{Parte I.}

\textbf{1.} $10^{12 - 7.46} = 10^{4.54} = \num{34700}$.

\textbf{2.} $N_\mathrm{He}/N_\mathrm{H} = 10^{10.914 - 12} = \num{0.0820}$, logo
$4N_\mathrm{He}/N_\mathrm{H} = \num{0.328}$, contra $Y/X =
0.2423/0.7438 = \num{0.326}$. Concordam até \qty{0.7}{\percent}; a
pequena diferença existe porque um átomo de hélio pesa um pouco menos do que quatro
átomos de hidrogênio.

\textbf{3.} $\log_{10}\bigl((1/350000)/\num{2.884e-5}\bigr) = -5.544 + 4.540 =
-1.00$.

\textbf{4.} $[\mathrm{O/H}] = -1.00 + 0.40 = -0.60$, logo
$\log\varepsilon(\mathrm{O}) = 8.69 - 0.60 = 8.09$. A estrela tem
$10^{8.09 - 6.46} = 42.7$ átomos de oxigênio por átomo de ferro, contra
$10^{1.23} = 17.0$ no Sol: $10^{0.4} = 2.5$ vezes mais.

\textbf{5.} $Z \approx 0.0139\times10^{-0.60} = \num{0.0035}$, um quarto do valor
solar.

\textbf{6.} O ferro tem um grande número de linhas no espectro de uma estrela fria, enquanto
o oxigênio tem poucas e fracas: o ferro é medido com muito mais
facilidade e precisão, e serve como representante dos metais.

\textbf{Parte II.}

\textbf{7.} Cada um se situa atrás de um isóbaro estável de $Z$ menor:
$^{122}\mathrm{Sn}$, $^{123}\mathrm{Sb}$ e $^{124}\mathrm{Sn}$ interrompem
as cadeias de decaimento do \omterm{def:b3:origin-of-elements:s-process}{processo r} antes que alcancem o telúrio. Os três situam-se
sobre o caminho s, que passa pelo telúrio de $^{122}\mathrm{Te}$ até
$^{126}\mathrm{Te}$.

\textbf{8.} \num{31.48}, \num{31.03} e \qty{30.75}{mb}; média
\qty{31.09}{mb}.

\textbf{9.} A média dos outros dois é $(31.48 + 30.75)/2 = \qty{31.11}{mb}$,
logo $N_{123} = 31.11/832 = \num{0.03740}$, apenas \qty{0.3}{\percent}
acima do valor medido, \num{0.0373}.

\textbf{10.} $1.329\times44.4 = \qty{59.0}{mb}$, quase o dobro do patamar.
$^{128}\mathrm{Te}$ é produzido apenas pelo \omterm{def:b3:origin-of-elements:s-process}{processo r}: o caminho s deixa o telúrio no
$^{127}\mathrm{Te}$, de modo que sua abundância foi formada pelo \omterm{def:b3:origin-of-elements:s-process}{processo r} e
não deve nada ao fluxo s.

\textbf{11.} $1 - 7.43/8.53 = \num{0.13}$: \qty{13}{\percent} do fluxo passou
ao redor de $^{148}\mathrm{Sm}$, por capturas nos instáveis
$^{147}\mathrm{Nd}$, $^{147}\mathrm{Pm}$ e $^{148}\mathrm{Pm}$ antes
que decaíssem até o samário.

\textbf{12.} Com uma camada fechada de 82 nêutrons, esses núcleos ligam um
nêutron extra fracamente e têm seções de choque de captura muito pequenas (\qty{4.0}{mb}
para $^{138}\mathrm{Ba}$). O fluxo se acumula sobre eles e apenas parte dele
prossegue: são gargalos.

\textbf{Parte III.}

\textbf{13.} $\mu_1 = 13.7/(13.7 + 33.4) = 13.7/47.1 = \num{0.291}$.

\textbf{14.} $\ln(1/0.291) = \num{1.235}$, logo $y = 0.0154/1.235 = \num{0.0125}$,
\num{0.81} vezes $Z_\odot$.

\textbf{15.} Em $\mu = 0.9$, $\ln(1/0.9) = 0.1054$, logo $Z = 0.0125\times0.1054 =
\num{0.00131}$. Como $0.00131/0.0154 = 0.0851$, $[\mathrm{Fe/H}] =
-1.07$. Para $\mu = 0.5$, $Z = 0.0125\times0.693 = \num{0.00864}$ e
$[\mathrm{Fe/H}] = -0.25$.

\textbf{16.} $Z = 0.00154$, logo $\mu = \mathrm{e}^{-0.00154/0.01247} =
\mathrm{e}^{-0.1235} = \num{0.884}$.

\textbf{17.} $1 - 0.884 = \num{0.116}$: \qty{12}{\percent} da massa.

\textbf{18.} $y = 0.0154/\ln 10 = \num{0.0067}$, cerca de metade disso.

\textbf{Parte IV.}

\textbf{19.} A massa das estrelas com \omterm{def:b3:origin-of-elements:metallicity}{metalicidade} abaixo de $Z$ é
$M(1 - \mathrm{e}^{-Z/y})$, logo, por unidade de $Z$, ela é
$(M/y)\,\mathrm{e}^{-Z/y}$. Como $Z = Z_\odot10^{[\mathrm{Fe/H}]}$,
$\mathrm{d}Z = Z\ln 10\,\mathrm{d}[\mathrm{Fe/H}]$, e, por unidade de
$[\mathrm{Fe/H}]$, a distribuição é proporcional a
$Z\mathrm{e}^{-Z/y}$. Sua derivada,
$(1 - Z/y)\mathrm{e}^{-Z/y}$, anula-se em $Z = y$: o pico ocorre em
$[\mathrm{Fe/H}] = \log_{10}0.810 = -0.09$.

\textbf{20.} $\num{0.116}/(1 - 0.291) = 0.1162/0.7091 = \num{0.164}$.

\textbf{21.} $Z/y = 1.235\times10^{-0.5} = 0.3905$, $1 - \mathrm{e}^{-0.3905} =
0.3233$, e $0.3233/0.7091 = \num{0.456}$.

\textbf{22.} $6670\times0.1638 = 1093$, contra 26: quarenta vezes mais do que o observado.

\textbf{23.} $26/902 = \qty{2.9}{\percent}$, ainda mais de cinco vezes abaixo
do previsto pela caixa fechada. As estrelas do modelo são as de vida longa, isto é, as de no máximo
cerca de uma massa solar, as quais sobreviveram todas desde que se formaram; entre
as estrelas F mais massivas, as mais antigas, que também são as mais
pobres em metais, morreram, e a amostra completa as subestima.

\textbf{24.} $(1 - 0.1^{0.1})/(1 - 0.1) = 0.2057/0.9 = \num{0.23}$: pior.
Aumentar a fração de gás também não ajuda: quando $\mu_1$ tende a 1
a fração tende a $0.1$, ainda mais de três vezes a
\qty{2.9}{\percent} observada. Nenhuma fração de gás salva a caixa fechada.

\textbf{25.} A caixa fechada prevê que
\textbf{\qty[text-series-to-math=true]{16}{\percent}} das estrelas de vida longa
da vizinhança solar têm $[\mathrm{Fe/H}] < -1$, contra
\qty{3}{\percent} no máximo observados. O disco não era uma caixa fechada: ele
cresceu por uma lenta acreção de gás pobre em metais, de modo que pouco gás jamais ficou
aguardando em baixa \omterm{def:b3:origin-of-elements:metallicity}{metalicidade} para formar estrelas.
\end{solution}
